\problemname{Vim}

It would be hard to find anything more geeky than the oldschool Vim --- an
editor in which we can do things that we could only dream of in ``regular''
editors. That is, if we are willing to invest a week of learning and a month of
practice so that the unusual but powerful commands it offers become our ``muscle
memory''. Well, just like in other editors, in Vim we have a cursor, which is
always located on one of the characters. Initially, it is located on the first
character of the text, and just like in other editors, in Vim there is a
clipboard, which is intended for storing (and indirectly copying and moving)
pieces of text. The clipboard is initially empty.

In this task, we will assume that there is exactly one \texttt{-} (minus)
character in the editor at the beginning, and our goal is to finish with exactly
$n$ consecutive minus signs. We will use four commands to help us with this:
\begin{itemize}
  \item \texttt{h}: If the cursor is on the first character, then this command
        does nothing, otherwise it moves the cursor to the character on the
        left.
  \item \texttt{l}: If the cursor is on the last character, then this command
        does nothing, otherwise it moves the cursor to the character on the
        right.
  \item \texttt{Y}: The sequence of characters extending from the cursor
        position to the end of the text is copied to the clipboard,
        ``overwriting'' any previous contents of the clipboard. (If the cursor
        is on the first character of the text, the entire text will be copied to
        the clipboard.)
  \item \texttt{P}: A copy of the text stored in the clipboard is inserted
        before the character on which the cursor is located, and the cursor is
        moved to the last inserted character. The contents of the clipboard are
        not changed. If the clipboard is empty, nothing happens.
\end{itemize}

Write a program that reads the number $n$ and prints the minimum number of
commands needed to finish with exactly $n$ consecutive \texttt{-} characters in
Vim under the described conditions. The program should also print an example of
a sequence with the minimum number of commands.

\section*{Input}
The first line contains the number of test cases $t$. The next $t$ lines contain
test cases with the desired string length $n$.

The following limits hold:
\begin{itemize}
  \item $1 \leq t \leq 100$
  \item $1 \leq n \leq 10^7$
\end{itemize}

\section*{Output}
For each test case, print the required minimum number of commands and an example
of such a sequence of commands on its own line.

Any sequence of commands achieving the minimum is accepted.

\section*{Scoring}
Your solution will be tested on a set of test groups.
To earn points for a group, you must pass all test cases in that group.

If you print only the correct number of commands, but no example or an incorrect
example of a sequence of commands, you will receive half the points for that
subtask.

\noindent
\begin{tabular}{| l | l | p{12cm} |}
  \hline
  \textbf{Group} & \textbf{Points} & \textbf{Constraints} \\ \hline
  $1$ & $20$ & $n \leq 100$ \\ \hline
  $2$ & $8$  & $n \leq 1\,000$ \\ \hline
  $3$ & $18$ & $n \leq 10^4$ \\ \hline
  $4$ & $18$ & $n \leq 10^5$ \\ \hline
  $5$ & $18$ & $n \leq 10^6$ \\ \hline
  $6$ & $18$ & No additional constraints. \\ \hline
\end{tabular}

\section*{Explanation of Samples}
As we can see in the following table, in the first case, the sequence
\texttt{YPYPhPYPPP} gives the correct result. The \texttt{=} character in the
Screen column represents the \texttt{-} character on which the cursor is
located.

\noindent
\begin{tabular}{| r | c | l | l |}
  \hline
  \textbf{Step} & \textbf{Command} & \textbf{Screen} & \textbf{Clipboard} \\ \hline
  $0$  &            & \verb|=|                     & (empty)      \\ \hline
  $1$  & \texttt{Y} & \verb|=|                     & \verb|-|     \\ \hline
  $2$  & \texttt{P} & \verb|=-|                    & \verb|-|     \\ \hline
  $3$  & \texttt{Y} & \verb|=-|                    & \verb|--|    \\ \hline
  $4$  & \texttt{P} & \verb|-=--|                  & \verb|--|    \\ \hline
  $5$  & \texttt{h} & \verb|=---|                  & \verb|--|    \\ \hline
  $6$  & \texttt{P} & \verb|-=----|                & \verb|--|    \\ \hline
  $7$  & \texttt{Y} & \verb|-=----|                & \verb|-----| \\ \hline
  $8$  & \texttt{P} & \verb|-----=-----|           & \verb|-----| \\ \hline
  $9$  & \texttt{P} & \verb|---------=------|      & \verb|-----| \\ \hline
  $10$ & \texttt{P} & \verb|-------------=-------| & \verb|-----| \\ \hline
\end{tabular}
